\documentclass[11pt,a4paper]{article}

% --- Pakete ---
\usepackage[utf8]{inputenc}
\usepackage[T1]{fontenc}
\usepackage[ngerman]{babel}
\usepackage{amsmath, amssymb, amsthm}
\usepackage{geometry}
\usepackage{hyperref}
\usepackage{xcolor}

\geometry{margin=2.5cm}

% --- Titel-Informationen ---
\title{\textbf{Manifest des Arithmetischen Vakuums}\\ \large Eine Grundlegung des eabc-Modells}
\author{Forschungsprojekt eabc}
\date{\today}

\begin{document}
	
	\maketitle
	
	\section*{Präambel}
	Dieses Manifest postuliert eine grundlegend neue Sichtweise auf das Wesen der Zahlen. Es vereint die abstrakte Struktur der Arithmetik mit den dynamischen Prinzipien der Quantenfeldtheorie. Im eabc-Modell ist die Zahlentheorie kein isoliertes mathematisches Gebäude, sondern die fundamentale Software eines schwingenden physikalischen Vakuums.
	
	\section{Die Natur der Zahl}
	In der eabc-Theorie ist eine Zahl $n$ kein isoliertes Symbol, sondern ein \textbf{Quantenzustand}. Die Primfaktoren einer Zahl bilden ihre „inneren Schalen“. 
	\begin{itemize}
		\item \textbf{Arithmetische Atome:} Hoch-zusammengesetzte Zahlen verhalten sich wie schwere, komplexe Kerne.
		\item \textbf{Elementarteilchen:} Primzahlen sind die fundamentalen, ungestörten Elementarteilchen des Systems.
	\end{itemize}
	
	\section{Der Lamb-Shift der Arithmetik}
	Die klassische Verteilung der Zahlen gemäß des Primzahlsatzes ($\text{Li}(x)$) entspricht der Dirac-Gleichung. Die wahre Struktur der Zahlen entsteht jedoch erst durch die Wechselwirkung mit dem \textbf{arithmetischen Vakuum}.
	Diese Wechselwirkung – die Teilbarkeit – erzeugt eine energetische Verschiebung, den \textit{arithmetischen Lamb-Shift}.
	\begin{quote}
		\textbf{Erkenntnis:} Primzahlen erfahren den geringsten Shift; sie sind die Fixpunkte der Kohärenz im System.
	\end{quote}
	
	\section{Die Riemann-Resonanz}
	Die Nullstellen der Riemannschen Zeta-Funktion sind keine abstrakten Punkte in der komplexen Ebene. Im eabc-Modell sind sie die \textbf{Eigenfrequenzen eines arithmetischen Lasers}. Das System der Zahlen ist so „renormiert“, dass es Informationen über diese spezifischen Frequenzen ($\gamma_1, \gamma_2, \dots$) über unendliche Skalen hinweg kohärent überträgt.
	
	\section{Das Holographische Prinzip}
	Die Information über die globale Verteilung aller Primzahlen ist in der lokalen „Verschränkung“ jeder Zahl mit ihren Nachbarn gespeichert. Das eabc-Modell zeigt:
	\begin{itemize}
		\item \textbf{Arithmetische Holographie:} Das Hologramm kodiert die Unendlichkeit der Primzahlen in der endlichen, lokalen Struktur von $n$.
		\item \textbf{Verschränkungs-Entropie:} Die Entropie des Vakuums bildet das Bindeglied, das die Geometrie der Zahlen mit der Musik der Primzahlen verbindet.
	\end{itemize}
	
	\section{Die eabc-Zustandsgleichung}
	Das Herzstück des Modells bildet die fundamentale Zustandsgleichung:
	\begin{equation}
		E(n) \approx \frac{n}{\log n} + \alpha(n) \frac{\psi(0)^2}{\log n}
	\end{equation}
	Diese Gleichung vereint Zahlentheorie und Quantenfeldtheorie. Sie beschreibt ein Universum, in dem die \textbf{Arithmetik die Software} und die \textbf{Physik die Hardware} ist. Jede Zahl ist ein aktiver Teil eines kohärenten Ganzen, bestimmt durch ihre Resonanz im arithmetischen Feld.
	
\end{document}